\documentclass[11pt]{beamer}
\usetheme[
  workplace=fi,
]{MU}

\usepackage{hyperref}
\usepackage[utf8]{inputenc}
\usepackage[english]{babel}
\usepackage[T1]{fontenc}
\usepackage{amsmath}
\usepackage{amsfonts}
\usepackage{amssymb}
\usepackage{graphicx}
\usepackage{listings}
\author{Jan Tušil}
\title{Formal Semantics of C++ language? Why?}
%\setbeamercovered{transparent} 
%\setbeamertemplate{navigation symbols}{} 
%\logo{} 
\institute{FI MU} 
\date{11.4.2018} 
%\subject{}


\lstset {
	basicstyle=\ttfamily
}

\newcommand{\backupbegin}{
   \newcounter{finalframe}
   \setcounter{finalframe}{\value{framenumber}}
}
\newcommand{\backupend}{
   \setcounter{framenumber}{\value{finalframe}}
}


\begin{document}

%%%%%%%%%%%%%%%%%%%%%%%%%%%%%%%%%%%%%%%%%%%%%%%%%%%%%%

% Phase 1: Attention
% If I go first:
% Good morning everyone. It is a prleasant surprise that there are
% so many of us in such an early hour. Did you sleep well?
% I hope you did, so you will not fall asleap now.
% If I go second:
% Good morning everyone. I hope you are enjoing the presentations
% as much as I do. Now you are hopefuly woke up enough not to fall asleap
% during my talk.

% Phase 2: Benefit
% My name is Honza Tušil. I could have named the talk 'Formal semantics of C++ and why you should care'.
% 

%My name is Honza Tušil
% Dobrý den, já jsem Honza Tušil a rád bych vám představil svoji diplomovou práci, 
% ve které jsem se zabýval vytvářením spustitelné formální sémantiky pro jazyk C++.
\begin{frame}
\titlepage
\end{frame}


% Phase 3: Credibility (20s)
% As I have said, my name is Honza Tusil
% and I am currently a Ph.D. candidate at FI MU.
% My research interested include software verification,
% programming language semantics and programming languages
% in general. I am also tutoring some C and C++ related classes
% and in past I did some part-time C/C++ development on embedded systems for a local company.
\begin{frame}{About me}
\begin{columns}

\begin{column}{0.30\textwidth}
\includegraphics[width=1.0\linewidth]{img/me.jpg}
\end{column}
\begin{column}{0.7\textwidth}
\begin{itemize}
\pause \item research: SW verification, PL semantics
\pause \item teaching: programming in C (PB071) and C++ (PB161)
\pause \item past: C/C++ @ embedded systems (local company)
\pause \item \url{jan.tusil@mail.muni.cz}
\pause \item \url{https://www.fi.muni.cz/~xtusil/}
\end{itemize}
\end{column}
\end{columns}
\end{frame}

% In the following ten minutes I would like to talk about the following. First I am going to show the current situation 
% of C++ analysis/verification tooling, then I will explain how
% is my Master thesis to this, and finally I am going to
% show the future work I would like to do.

\begin{frame}{The structure}
\tableofcontents
\end{frame}

\AtBeginSection[]
{
  \begin{frame}
    \frametitle{Osnova}
    \tableofcontents[currentsection]
  \end{frame}
}

%\section{The Context}
\section{Why bother?}

% Modern programming languages are really powerful - but also
% really complex. The specification of the C++ language alone
% has circa 380 pages, and the language and its standard library
% together has 1400 pages in total.
% As the time goes, some features are being added to the language,
% other are being removed. For the C++ language there is
% a nice map of the features. It looks like this.
\begin{frame}{The C++ lands}
\pause
\includegraphics[trim=100 900 1700 300, clip, width=1.0\linewidth]{img/cppmap.png}
%TODO link
\end{frame}

% And this is the full map
\begin{frame}{The C++ lands (full)}
\includegraphics[width=1.0\linewidth]{img/cppmap.png}
\end{frame}


% 2 min


% The C++ language is complex and we
% - the programmers - make mistakes.
% Therefore we have a plenty of tools which help as to find
% bugs in our programs, or to prove their absence from the code.
% The last two tools in the list are being actively developed
% here on Faculty on Informatics.

% The tools can give programmers a valuable information.
% For example, that their program contains an error.
% Or that some code is suspicious, and it could be an error.
% Or they can say that the program is error-free.
\begin{frame}{The complexity of C/C++}
\begin{columns}

\begin{column}{0.5\textwidth}
Some of the tools:
\begin{itemize}
\pause \item clang-tidy
\pause \item frama-c
\pause \item cbmc
\pause \item divine
\pause \item symbiotic
\end{itemize}
\end{column}

\begin{column}{0.5\textwidth}
\pause
Possible outputs:
\begin{itemize}
\pause \item error
\pause \item suspicious code, possible error
\pause \item correct
\pause \item unknown
\end{itemize}
\end{column}

\end{columns}
\end{frame}

% This, however, leads to some important questions.
% Are they correctly implemented?
% Do these tools understand the C++ language correctly?

% There is a competetion of verification tools, called SV-COMP.
% In this competetion, there are two set of C programs:
% a set of correct programs, and a set of incorrect programs.
% However, two years ago it was discovered that some of the programs
% in the 'correct' set are actually incorrect (they contain
% an undefined behaviour). So the tools that classify such programs
% as correct gain some points in the competetion, but actually
% they contain an error - they do not understand the C language
% correctly.
\begin{frame}{Some disturbing questions}

\begin{itemize}
\pause \item Are they correctly implemented?
\pause \item Do they really ,,understand'' C++?
\pause \item SV-COMP: programs containing undefined behaviour in the set of correct ones\footnote[frame]{https://runtimeverification.com/blog/?p=200}
\end{itemize}

\end{frame}


% Jeden z možných přístupů, jak mohou nástroje rozumět programovacímu jazyku,
% je delegovat toto porozumění na překladač. Překladač přeloží program v C/C++
% do nízkoúrovňového jazyka (třeba do assembleru) a nástroj už musí rozumět
% jen tomuto nízkoúrovňovému jazyku. V posledních letech je hodně populární
% platforma překladačů llvm spolu s jejím C/C++ frontendem, který se jmenuje Clang.
% Tento přístup volí například Divine, Klee, a Symbiotic.

\begin{frame}{Solution no. 1 - use a compiler}
Use a translator from the desired language to some simpler language.

E.g. a compiler (clang/llvm). This way the tool does not have to understand the complex C++ language, because the compiler understands it.

\end{frame}



% Osobně vidím dvě nevýhody tohoto přístupu: za prvé, nelze takto ověřit korektnost programu
% vzhledem k programovacímu jazyku; a za druhé, verifikace některých vysokoúrovňových vlastností
% a konstrukcí se stává téměř, ne-li úplně, nemožná.


\begin{frame}{Limitations}

\begin{itemize}
\pause \item takto ale nelze ověřit korektnost programu vzhledem k jazyku
\pause \item jak verifikovat šablony?
\pause \item jiný přístup: samostatná formální sémantika
\end{itemize}


\end{frame}

% Jiným přístupem je mít samostatnou, spustitelnou formální sémantiku daného jazyka.

\begin{frame}{Solution no. 2 - have a formal semantics}

\end{frame}


% Právě k popisu formální sémantiky programovacího jazyka slouží
% framework zvaný "K". Tento framework umožňuje dát jazyku
% small-step operační sémantiku, založenou na přepisování termů.
% Nástroje jako parser, interpreter, či model checker
% je možné pomocí tohoto frameworku naprogramovat obecně, nezávisle
% na syntaxi a sémantice konkrétního jazyka.
% Zrovna tyto zmiňované nástroje jsou již naprogramovány a je možné je téměř automaticky instanciovat
% pro libovolný jazyk, pro který je v K napsaná sémantika.
% 

\begin{frame}{K Framework}
\includegraphics[width=1.0\linewidth]{img/kidea.png}
\end{frame}

% V K již byla popsána sémantika Pythonu, JavaScriptu, C, a dalších jazyků.
% Nad sémantikou jazyka C byl postaven komerční nástroj pro detekci
% nedefinovaného chování RV-Match a tato sémantika je poslední
% dva roky rozšiřována o podporu jazyka C++, které je zatím jen v experimentální fázi.
\begin{frame}{C/C++ @ K}
\begin{itemize}
\item Already formalized: Python, JavaScript, C
\pause \item Application of the C semantics: RV-Match \\
\includegraphics[width=0.3\linewidth]{img/rvmatch.png}
\pause \item current work: extend the C semantics to C++
\end{itemize}
\end{frame}

% Zadání, cíle, konkrétní věci, nalezené chyby apod.
\section{Ms Thesis - The Results}


% Cílem mé práce bylo rozšířit stávající prototyp
% C++ sémantiky v K o dosud neimplementované vlastnosti tohoto jazyka.

% První vlastností, kterou jsem zamýšlel implementovat,
% jsou výčtové typy. V projektu sice výčtové typy byly implementovány,
% ale jen pro jazyk C. V jazyce C++ se výčtové typy chovají trošku jinak,
% navíc ve standardu C++11 (z roku 2011) přibyla nová varianta těchto typů.


\begin{frame}[fragile]{Výčtové typy (\texttt{enum})}
\begin{lstlisting}[language=C++]
enum E0 { A = 5, B = A + 3 };
enum E1 : char;
enum class E2 : unsigned int;
enum class E3;
enum E5 { A, B } __attribute__((packed));
std::underlying_type_t<E5> x = 5;
\end{lstlisting}
\end{frame}



% Další vlastností, kterou jsem chtěl implementovat, byla podpora vykonávání kódu
% v době překladu. Od standardu C++11 je možné označit některou z funkcí
% jako 'constexpr', což v praxi znamená, že volání takovéto funkce je možné
% vyhodnotit již v době překladu. Výpočet pak není nutné dělat za běhu,
% což může výrazně pomoci jak výkonu programu, tak čitelnosti jeho kódu.
% Výsledek výpočtu je navíc možné použít i na místech, kde jazyk vyžaduje použití konstanty.
% A ano - jazyk C++ je díky této vlastnosti turingovsky úplný už v době překladu.
% Jenže on byl turingovsky úplný již dávno předtím.

\begin{frame}[fragile]{Vykonávání kódu v době překladu (\texttt{constexpr})}
\begin{lstlisting}
constexpr int fact(int n) {
	return n <= 1 ? 1 : n * fact(n - 1);
}
int main() {
	constexpr int f_4 = fact(4);
	enum class Facts {
		F_4 = f_4,
		F_5 = fact(5)
	};
	assert(int(Facts::F_4) == 24);
	assert(int(Facts::F_5) == 120);
}
\end{lstlisting}
\end{frame}


% Přestože jsem to původně neměl v úmyslu, na popud jednoho z hlavních vývojářů projektu
% jsem implementoval vlastnost, která se nazývá "nulová inializace tříd".
% Myslel jsem si, že to bude hračka, ale zabralo mi to dost dlouhou dobu.
% Jedním z důvodů byla i chyba v projektu, která se projevovala až tehdy,
% když jsem se pokoušel nulovou inicializaci naimplementovat. Chybou byl
% nedeterminismus v sémantice, který umožňoval interpretovat některé programy vícero způsoby.

\begin{frame}[fragile]{Nulová inicializace tříd}
\begin{lstlisting}
struct S {
	virtual ~S() = default;
	int x;
};

S s1;
S s2{};
assert(s1.x == 0); // ok
assert(s2.x == 0); // fail
\end{lstlisting}
\end{frame}

% Tím se dostáváme k poslednímu bodu - opravě některých stávajících chyb.

% Projekt C++ sémantiky v K je stále v experimentální fázi,
% a tak se při práci častokrát stalo, že jsem narázil
% na nějakou chybu či nedokončenou implementaci. Opravoval jsem zejména
% ty, které mi nějak překážely.

% Ještě doplním, že z těchto čtyř bodů byly jen první dva
% součástí mého oficiální zadání.
\begin{frame}{Oprava stávajících chyb v projektu}
\includegraphics[width=0.5\linewidth]{img/bug.jpg}
\end{frame}


% Jak to dopadlo? Enumy jsou kompletně implementovány,
% většina implementace je již nějakou dobu začleněna v upstreamu.
% Vyhodnocování funkcí v době překladu - aka constexpr -
% je dost komplexní vlastnost, a její hezká implementace
% vyžaduje hluboké zásahy do projektu. Vytvořil jsem prototyp,
% který implementuje ty koncepčně náročné aspekty této vlastnosti (včetně rekurze apod).
% Není ale zatím připravený pro začlenění.
% Nulová inicializace je hotova, nicméně před začleněním do upstreamu
% musí být ještě opravana jedna s ní související chyba.
% Když jsme u chyb - pár jsem jich opravil, některé jsou již začleněny,
% jiné mám v úmyslu začlenit do upstreamu později.
% Jako bonus jsem objevil nějaké chyby v jiných projektech
% - jednu v překladači GCC, jednu v Clang a jednu v samotném standardu jazyka C++.
% Chyby u překladačů jsem pochopilteně nahlásil, s chybou ve standardu to je trochu složitější.
\begin{frame}{Výsledky}
\begin{itemize}
\pause \item \texttt{enum} - vše hotovo, většina je v upstream
\pause \item \texttt{constexpr} - elegantní proof of concept
\pause \item nulová inicializace - hotovo
\pause \item opravy chyb - některé přijaty do upstreamu
\pause \item nalezené chyby v GCC, Clang, standardu C++
\end{itemize}
\end{frame}

\section{Future work}


% Co teď? Je třeba začlenit hotové věci do upstreamu,
% dokončit rozdělané vlastnosti a implementovat všechny ostatní,
% které ještě nejsou - a je jich dost. Je také třeba vymyslet,
% jak budeme v projektu řešit vývoj jazyka - jestli budeme třeba podporovat
% nějaké přepínače pro výběr standardu, nebo raději pojedeme vždy podle nejnovější verze.
% Mám v úmyslu ve vývoji sémantiky pokračovat, jsem dohodnutý s firmou RuntimeVerification inc
% (která má nad C sémantikou postavený onen nástroj RV-match, který jsem zmiňoval
% na začátku prezentace) na tom, že vývoj finančně podpoří.
%
% Kromě toho bych rád implementoval podporu kontraktů a jejich verifikace
% (jak statické, tak dynamické). Verifikace kontraktů
% u programů s vysokou úrovní abstrakce je vlastně ten důvod,
% kvůli kterému jsem na sémantice jazyka C++ začal pracovat.
\begin{frame}{Future of the semantics}
\begin{itemize}
\pause \item Merge the rest to upstream.
\pause \item Finalize \texttt{constexpr}.
\pause \item Templates, threads, bit fields, fold expressions, \ldots
\pause \item How to deal with language changes (C++14, 17, 20, \ldots) ?
\pause \item An official contract with RuntimeVerification inc.
\end{itemize}
\end{frame}

% TODO maybe: correctness of an implementation of an analysis

% Rád bych tedy ve zbývajících dvou minutách ukázal,
% jak si takové kontrakty představuji. Na slajdu vidíme
% deklaraci šablonové funkce sort(), která bere dva iterátory
% libovolného typu a seřadí rozsah jimi zadaný.
% Tato funkce požaduje, aby iterátory tvořily platný rozsah,
% tedy aby bylo možné doiterovat od prvního ke druhému.
% Funkce garantuje, že po jejím dokončení bude tento rozsah seřazený.
% Jméno is_sorted v kontraktu této funkce je název úplně obyčejné
% šablonované funkce v jazyce C++, která zkontroluje,
% že rozsah jí předaný je seřazený.
%
% Kontrakty by tedy měly umět pracovat s běžnými výrazovými
% prostředky jazyka C++, zejména s šablonami, funkcemi,
% iterátory a jinými uživatelsky definovanými abstrakcemi.
% Jak na to - uvidíme.
\begin{frame}[fragile=singleslide]{Contracts on generics (C++ templates)}
\begin{lstlisting}
template <typename It>
bool is_sorted(It begin, It end);

template <typename It>
bool valid_range(It begin, It end);

template <typename It>
void sort(It begin, It end)
  [[ expects axiom: valid_range(begin, end) ]]
  [[ ensures: is_sorted(begin, end) ]] ;
\end{lstlisting}
\pause This is the actual syntax proposed for C++20.
\end{frame}


\begin{frame}{Kontrakty}
\begin{itemize}
\pause \item Jak jim dát sémantiku?
\pause \item Volání běžných funkcí v pre/postcondition.
\pause \item A co když mají efekty? A co když ty jsou ohraničené?
\pause \item Jak je verifikovat?
\pause \item Jak využít existující abstrakce? A vzory?
\end{itemize}
\end{frame}


\begin{frame}{Concepts and Contracts in C++20}

\end{frame}

\begin{frame}{Zdroje}
\begin{itemize}
\item Mapa C++ \url{https://fearlesscoder.blogspot.cz/2017/02/the-c17-lands.html}
K diagram \url{http://fsl.cs.illinois.edu/index.php/From\_Rewriting\_Logic,\_to\_Programming\_Language\_Semantics,\_to\_Program\_Verification} \\
\item RV-Match logo \url{https://runtimeverification.com/match/}
\item Brouk \url{http://cliparting.com/free-bug-clipart-30526/}
\end{itemize}
\end{frame}

\backupbegin

\section{Bonus for C++ folks}


% Tohle bych mohl mozna presunout do bonusu.
% Tady bych jenom dodal, že nalezení chyb v překladačíh
% nebo standardech programovacích jazyků není až tak moc neobvyklé,
% zvlášť v kontextu formalizace těchto jazyků - JavaScript, Ethereum.
%\subsection{Chyby nalezené v jiných projektech}

% Během práce jsem hodně pracoval se standardem jazyka C++,
% a když jsem o C++ zjišťoval nové věci, prakticky jsem je zkoušel
% na malých příkladech. Kladl jsem si otázku, "co se stane, když..."?
% Tímto způsobem jsem našel chybu v překladači GCC.
% Kód na slajdu je platný C++ kód, a Clang nebo MSVC jej bez problémů zkompilují,
% nicméně GCC jej odmítne zparsovat. Chybu jsem nahlásil do bugzilly projektu
% a jeden z hlavních vývojářů ji potvrdil.
\begin{frame}[fragile=singleslide]{GCC bug}
\begin{lstlisting}
enum E : int;
enum ::E : int{};
\end{lstlisting}
\end{frame}

% Další chybu jsem objevil v překladači clang.
% Chyba se týká toho, v jakém scope žijí konstanty výčtového typu
% (v našem případě konstanty A a B).
%
% Podle jednoho z vývojářů překladače clang, konkrétně správce
% C++ frontendu tohoto překladače, se příčina této chyby projevuje
% i v jiných případech, ve kterých nemá vliv na korektnost, ale jen
% na výkon překladače.
\begin{frame}[fragile=singleslide]{Clang bug}
\begin{lstlisting}
namespace N { enum E : int; }
enum N::E : int {A,B};
namespace N {
    int x = int(::N::B); // 1
}
int y = int(::A); // 2
int z = int(A); // 3
\end{lstlisting}
\end{frame}

\begin{frame}[fragile=singleslide]{C++ standard bug}
\begin{lstlisting}
enum class E {
    A,
    B = E::A // 1
};
\end{lstlisting}
Jaký typ má \lstinline|E::A| v (1)?
\end{frame}

% Btw v C++ sémantice stále ještě není implementován switch.
% Po otázkách - kdyby chtěli vědět něco dalšího, rád si s nima popovídám.
% Co jsem se naučil: komunikace na dálku je fuška.
% S lidma z RV a UIUC se mi spolupracovalo dobře.

\backupend

\end{document}